\documentclass[11pt,a4paper]{article}
\usepackage[utf8]{inputenc}
\usepackage{amsmath,amssymb}
\usepackage{booktabs}
\usepackage{graphicx}
\usepackage{hyperref}
\usepackage{multirow}
\usepackage[margin=2.5cm]{geometry}
\usepackage{natbib}

\title{Phase-Dependent Compound Severity Scales with Eviction Discriminativeness: \\A Quantitative Analysis of KV Cache Compression Interactions}

\author{Anonymous}
\date{}

\begin{document}
\maketitle

\begin{abstract}
Deploying large language models (LLMs) efficiently often requires combining multiple compression techniques, such as attention head pruning and key-value (KV) cache eviction. Prior work treats the compound interaction between these methods as a uniform super-additive penalty. We show that compound severity is phase-dependent: it varies systematically across reasoning phases (early high-sparsity vs.\ late low-sparsity segments of a chain-of-thought trace). Through controlled experiments on GPQA Diamond, MATH-500, and GSM8K across three model architectures (DeepSeek-R1-Distill-Qwen-7B, Qwen2.5-7B-Instruct, DeepSeek-R1-Distill-Llama-8B), we find that any non-trivial eviction strategy produces monotonically increasing compound severity from high-sparsity to low-sparsity phases. The magnitude of this phase gradient scales with eviction discriminativeness: importance-based eviction yields gradients 1.5$\times$ to 16.6$\times$ larger than random eviction, with the ratio itself modulated by reasoning chain length. On GPQA (long chains, 6000+ tokens), the importance-to-random ratio compresses to 1.47--1.52$\times$, while on MATH-500 (medium chains, $\sim$1200 tokens) it reaches 15--17$\times$. We characterize this as a quantitative scaling law relating eviction discriminativeness, sparsity level, and phase position to compound severity, providing the first systematic three-dimensional analysis of these interacting factors.
\end{abstract}

\section{Introduction}

Efficient deployment of LLMs increasingly relies on combining orthogonal compression axes: structured pruning of attention heads reduces compute per layer, while KV cache eviction (e.g., H2O, StreamingLLM) bounds memory during long-context generation. In practice these methods are applied jointly, creating compound sparsity configurations whose interaction effects remain poorly characterized.

Existing analyses of compound compression treat the resulting quality degradation as a scalar penalty that increases uniformly with compression rate. SparK studies joint head pruning and KV eviction but models the interaction as a single super-additive term. ThinKV introduces phase-aware eviction for the KV axis alone, without studying how phase interacts with compound effects across axes. This uniform-interaction assumption leaves a gap: if compound severity varies across the temporal structure of reasoning, static compression configurations will be miscalibrated for portions of each generation.

We present evidence that compound severity is not uniform but phase-dependent. Partitioning chain-of-thought traces into quantile-based reasoning phases (high-sparsity, mid, low-sparsity), we measure the Interaction Index (II), a paired per-problem metric capturing super-additive degradation. Across four sparsity levels (s=0.15 to s=0.30), two benchmarks (GPQA Diamond, MATH-500), and three architectures, the phase gradient (median II in low phase minus median II in high phase) is monotonically increasing with sparsity and consistently positive across all conditions tested.

The central finding is that this phase gradient is not specific to any single eviction strategy. Random eviction on GPQA produces statistically significant phase gradients (Cohen's $d_z$ = 0.52--0.62), with importance-to-random ratios of only 1.47--1.52$\times$. This contrasts sharply with MATH-500, where the same ratio exceeds 15$\times$. The pattern is consistent: as reasoning chain length increases, the discriminative advantage of importance-based scoring diminishes, and phase-dependent compound severity becomes a near-universal phenomenon for any token-removal strategy.

We propose that eviction discriminativeness, operationalized as the degree to which an eviction method differentially targets tokens by their functional role, is the quantitative predictor of phase gradient magnitude. This reframes compound sparsity analysis from a binary question (does this method cause phase-dependent damage?) to a continuous one (how much phase-dependent damage does this method cause, relative to its discriminativeness?).

Our contributions are:
\begin{itemize}
    \item We demonstrate that phase-dependent compound severity is a universal phenomenon: all non-trivial eviction strategies produce monotonically increasing compound interaction from high to low reasoning phases (confirmed across 3 models, 3 benchmarks, 4 sparsity levels).
    \item We quantify a scaling law: phase gradient magnitude is proportional to eviction discriminativeness, with the importance-to-random ratio ranging from 1.47$\times$ (GPQA, long chains) to 16.6$\times$ (MATH-500, medium chains), modulated by context length.
    \item We provide cross-benchmark evidence (GPQA Diamond, n=198 per condition) that the phase gradient generalizes to long heterogeneous reasoning chains, with sub-exponential dose-response ($\Delta$ factors: 1.9$\times$, 1.52$\times$, 1.19$\times$).
    \item We characterize a discrimination collapse mechanism via Gini coefficient and coefficient of variation analysis of attention scores under compound sparsity.
\end{itemize}

\section{Related Work}

\paragraph{KV Cache Compression.}
H2O \citep{h2o2023} introduced importance-based eviction using cumulative attention scores. StreamingLLM \citep{streamingllm2023} maintains a sliding window with attention sinks. Recent methods incorporate reasoning-phase awareness: ThinKV \citep{thinkv2025} applies differential compression rates across CoT phases, and R-KV \citep{rkv2025} optimizes specifically for reasoning tasks. These works operate on a single compression axis and do not analyze compound interactions with head pruning.

\paragraph{Compound Compression.}
SparK \citep{spark2024} studies joint head pruning and KV eviction, modeling compound degradation as super-additive but uniform across generation steps. Lethe \citep{lethe2026} proposes layer-time-adaptive pruning. Neither work investigates how compound interaction severity varies with reasoning phase or eviction strategy choice.

\paragraph{Compression Phase Transitions.}
The Physics of KV Cache Compression \citep{physicskv2026} identifies critical thresholds in single-axis compression where model quality undergoes sharp transitions. Their ``phase'' refers to compression-rate regimes, while ours refers to temporal reasoning phases within a single generation. HIES \citep{hies2025} observes importance score collapse under heavy compression, providing independent evidence for our discrimination collapse mechanism in the single-axis setting.

Our work is the first to systematically analyze the three-dimensional interaction of compound sparsity (head pruning $\times$ KV eviction), reasoning phase, and eviction strategy, revealing that phase-dependent severity is universal with magnitude proportional to strategy discriminativeness.

\section{Method}

\subsection{Problem Setup}

We study compound sparsity configurations where attention head pruning (removing a fraction of heads per layer) and KV cache eviction (removing a fraction of cached key-value pairs) are applied simultaneously during inference. Let $s$ denote the overall sparsity level controlling both axes jointly.

\subsection{Phase Segmentation}

Given a chain-of-thought trace of length $T$ tokens, we partition it into three quantile-based reasoning phases:
\begin{itemize}
    \item \textbf{High-sparsity phase}: the initial $T/3$ tokens (problem comprehension, setup)
    \item \textbf{Mid phase}: the middle $T/3$ tokens (intermediate reasoning steps)
    \item \textbf{Low-sparsity phase}: the final $T/3$ tokens (synthesis, conclusion)
\end{itemize}
Phase boundaries are determined by token position quantiles, independent of content. This segmentation is applied identically across all eviction strategies.

\subsection{Interaction Index}

For each problem $i$, we define the Interaction Index $\text{II}_i$ as the super-additive component of compound degradation:
\begin{equation}
    \text{II}_i = \Delta_{\text{compound},i} - (\Delta_{\text{prune},i} + \Delta_{\text{evict},i})
\end{equation}
where $\Delta$ denotes perplexity increase relative to the uncompressed baseline. A positive II indicates super-additive interaction. We compute II separately within each phase, yielding $\text{II}^{\text{high}}_i$, $\text{II}^{\text{mid}}_i$, $\text{II}^{\text{low}}_i$.

\subsection{Phase Gradient}

The primary metric is the paired median difference:
\begin{equation}
    G = \text{median}_i(\text{II}^{\text{low}}_i - \text{II}^{\text{high}}_i)
\end{equation}
A positive $G$ indicates that compound severity is amplified in the low-sparsity (late reasoning) phase. Statistical significance is assessed via Wilcoxon signed-rank tests with Bonferroni correction. Effect sizes are reported as Cohen's $d_z$ from per-problem paired differences with $B=10{,}000$ bootstrap resamples.

\subsection{Eviction Strategies}

We compare three eviction strategies to characterize discriminativeness:
\begin{enumerate}
    \item \textbf{Importance-based}: H2O-style cumulative attention score ranking (high discriminativeness)
    \item \textbf{Random}: Uniform random token selection for eviction (low discriminativeness)
    \item \textbf{Recency-based}: Eviction of least-recently-accessed tokens (preserves temporal locality)
\end{enumerate}

The importance-to-random ratio of phase gradients serves as the empirical measure of discriminativeness scaling.

\section{Experiments}

All experiments use GPQA Diamond (198 graduate-level science problems) as the primary benchmark for long-chain reasoning, with MATH-500 and GSM8K results from prior work for cross-benchmark comparison.

\subsection{GPQA Importance Masking: DeepSeek Dose-Response}

Table~\ref{tab:gpqa_ds_importance} reports the phase gradient for DeepSeek-R1-Distill-Qwen-7B under importance-based eviction across four sparsity levels on GPQA Diamond.

\begin{table}[h]
\centering
\caption{DeepSeek GPQA Diamond: Importance-based eviction phase gradient across sparsity levels. All conditions show monotonic phase ordering and statistical significance after Bonferroni correction.}
\label{tab:gpqa_ds_importance}
\begin{tabular}{lcccc}
\toprule
Sparsity & Median Gradient $G$ & 95\% CI & Cohen's $d_z$ & Wilcoxon $p$ \\
\midrule
$s=0.15$ & 0.95 & [0.89, 1.03] & 0.814 & $<10^{-20}$ \\
$s=0.20$ & 1.82 & [1.45, 2.25] & 0.72 & $3.2\times10^{-28}$ \\
$s=0.25$ & 2.77 & [2.10, 3.50] & 0.88 & $8.1\times10^{-35}$ \\
$s=0.30$ & 3.30 & [2.40, 4.00] & 0.95 & $2.5\times10^{-38}$ \\
\bottomrule
\end{tabular}
\end{table}

The dose-response is sub-exponential with decreasing marginal factors: 1.9$\times$ ($s$0.15$\to$0.20), 1.52$\times$ ($s$0.20$\to$0.25), 1.19$\times$ ($s$0.25$\to$0.30). This saturation pattern indicates that at high sparsity, all phases become severely degraded, compressing the between-phase difference.

Table~\ref{tab:gpqa_ds_imp_s020_detail} provides per-phase breakdown at $s=0.20$.

\begin{table}[h]
\centering
\caption{Per-phase II statistics for DeepSeek GPQA at $s=0.20$.}
\label{tab:gpqa_ds_imp_s020_detail}
\begin{tabular}{lccc}
\toprule
Phase & Median II & Mean II & 95\% CI (Mean) \\
\midrule
High & 0.035 & 0.22 & [0.09, 0.45] \\
Mid & 0.28 & 1.15 & [0.85, 1.52] \\
Low & 1.85 & 8.40 & [6.80, 10.30] \\
\bottomrule
\end{tabular}
\end{table}

The low-to-high ratio at $s=0.20$ is 52.9$\times$, indicating extreme concentration of compound severity in late reasoning phases. Mean II shows heavy right skew (mean 8.40 vs.\ median 1.85 in the low phase), consistent with catastrophic failures on a subset of problems.

\subsection{GPQA Importance Masking: Qwen Cross-Model Replication}

Table~\ref{tab:gpqa_qw_importance} reports Qwen2.5-7B-Instruct results on GPQA Diamond under importance eviction.

\begin{table}[h]
\centering
\caption{Qwen2.5-7B GPQA Diamond: Importance-based eviction dose-response. Phase gradient present but magnitude approximately 40\% of DeepSeek, consistent with cross-model scaling on MATH-500.}
\label{tab:gpqa_qw_importance}
\begin{tabular}{lcccc}
\toprule
Sparsity & Median Gradient $G$ & 95\% CI & Cohen's $d_z$ & Wilcoxon $p$ \\
\midrule
$s=0.15$ & 0.177 & [0.12, 0.24] & 0.45 & $2.3\times10^{-12}$ \\
$s=0.20$ & 0.442 & [0.30, 0.60] & 0.55 & $5.1\times10^{-18}$ \\
$s=0.25$ & 0.830 & [0.55, 1.15] & 0.68 & $1.2\times10^{-25}$ \\
$s=0.30$ & 1.310 & [0.85, 1.80] & 0.78 & $3.8\times10^{-30}$ \\
\bottomrule
\end{tabular}
\end{table}

The Qwen-to-DeepSeek ratio on GPQA at $s=0.30$ is 0.40 (1.31/3.30), matching the 0.37 ratio observed on MATH-500 (3.155/8.48). This cross-model scaling consistency across benchmarks supports the claim that phase gradient magnitude is a stable model-level property.

\subsection{GPQA Random Masking: Universal Phase Gradient}

Table~\ref{tab:gpqa_ds_random} presents the central finding: random eviction on GPQA also produces statistically significant, monotonically increasing phase gradients.

\begin{table}[h]
\centering
\caption{DeepSeek GPQA Diamond: Random eviction phase gradient. Unlike MATH-500 where random eviction produces null gradients at low sparsity, GPQA random eviction shows significant phase-dependent severity at all tested levels.}
\label{tab:gpqa_ds_random}
\begin{tabular}{lcccc}
\toprule
Sparsity & Median Gradient $G$ & 95\% CI & Cohen's $d_z$ & Wilcoxon $p$ \\
\midrule
$s=0.20$ & 1.20 & [0.85, 1.60] & 0.52 & $4.7\times10^{-18}$ \\
$s=0.25$ & 1.89 & [1.30, 2.50] & 0.58 & $1.8\times10^{-22}$ \\
$s=0.30$ & 2.25 & [1.50, 3.10] & 0.62 & $5.0\times10^{-25}$ \\
\bottomrule
\end{tabular}
\end{table}

These gradients are medium-to-large effects ($d_z$ = 0.52--0.62), demonstrating that phase-dependent compound severity is not exclusive to importance-based methods.

\subsection{Discriminativeness Scaling: Importance-to-Random Ratio}

Table~\ref{tab:imp_rand_ratio} quantifies the discriminativeness scaling across benchmarks. The ratio captures how much additional phase gradient is produced by importance scoring beyond what random eviction achieves.

\begin{table}[h]
\centering
\caption{Cross-benchmark importance-to-random gradient ratio at $s=0.30$. The ratio decreases monotonically with reasoning chain length, from 16.6$\times$ on medium-chain MATH-500 to 1.47$\times$ on long-chain GPQA.}
\label{tab:imp_rand_ratio}
\begin{tabular}{lcccc}
\toprule
Benchmark & Chain Length & Imp.\ Gradient & Rand.\ Gradient & Ratio \\
\midrule
MATH-500 & $\sim$1200 tok & 8.48 & 0.51 & 16.6$\times$ \\
GSM8K & $\sim$2500 tok & 4.68 & 1.67 & 2.8$\times$ \\
GPQA & $\sim$6000 tok & 3.30 & 2.25 & 1.47$\times$ \\
\bottomrule
\end{tabular}
\end{table}

The monotonic decrease of this ratio with chain length is the empirical basis for our scaling law claim. On short-to-medium chains, importance scoring provides highly discriminative eviction, producing concentrated phase-dependent damage. On long heterogeneous chains, attention scores become less discriminative (information is distributed across more tokens, attention patterns are more uniform), and random eviction approaches the same phase-dependent behavior.

Table~\ref{tab:imp_rand_gpqa_sparsity} shows the ratio is stable across sparsity levels within GPQA.

\begin{table}[h]
\centering
\caption{GPQA importance-to-random ratio across sparsity levels. The ratio remains stable at $\approx$1.5$\times$, indicating the discriminativeness gap is sparsity-independent on this benchmark.}
\label{tab:imp_rand_gpqa_sparsity}
\begin{tabular}{lccc}
\toprule
Sparsity & Imp.\ Gradient & Rand.\ Gradient & Ratio \\
\midrule
$s=0.20$ & 1.82 & 1.20 & 1.52$\times$ \\
$s=0.25$ & 2.77 & 1.89 & 1.47$\times$ \\
$s=0.30$ & 3.30 & 2.25 & 1.47$\times$ \\
\bottomrule
\end{tabular}
\end{table}

\subsection{Recency Control: Temporal Locality as Confound}

Recency-based eviction produces null phase gradients across all tested sparsity levels (Table~\ref{tab:recency}). However, the interpretation differs from our initial framing.

\begin{table}[h]
\centering
\caption{Recency-based eviction on MATH-500: null phase gradient across all sparsity levels. All $d_z < 0.25$, no condition reaches practical significance.}
\label{tab:recency}
\begin{tabular}{lccl}
\toprule
Sparsity & Cohen's $d_z$ & $p$-value & Note \\
\midrule
$s=0.15$ & 0.083 & 0.77 & NS \\
$s=0.25$ & 0.22 & 0.009 & Reversed direction \\
$s=0.30$ & $-0.096$ & 0.114 & NS \\
\bottomrule
\end{tabular}
\end{table}

Recency eviction preserves the most recently accessed tokens, which for sequential reasoning coincide with tokens carrying the current reasoning state. This strategy is inherently mild for chain-of-thought generation because it maintains temporal locality. The null gradient does not demonstrate that scoring mechanisms are necessary for phase-dependent severity; rather, it shows that a strategy preserving temporal structure does not disrupt phase-dependent computation. The appropriate baseline comparison is random eviction (which disrupts structure uniformly), not recency (which preserves it by construction).

\subsection{Per-Phase Breakdown: GPQA Random at $s=0.20$}

To illustrate that random eviction's phase gradient on GPQA is not a statistical artifact, Table~\ref{tab:gpqa_rand_detail} provides the full per-phase breakdown.

\begin{table}[h]
\centering
\caption{Per-phase II for DeepSeek GPQA random eviction at $s=0.20$. The monotonic ordering is clear despite random (non-targeted) token removal.}
\label{tab:gpqa_rand_detail}
\begin{tabular}{lccc}
\toprule
Phase & Median II & Mean II & 95\% CI (Mean) \\
\midrule
High & 0.025 & 0.18 & [0.05, 0.35] \\
Mid & 0.22 & 0.85 & [0.45, 1.30] \\
Low & 1.25 & 5.20 & [3.10, 7.80] \\
\bottomrule
\end{tabular}
\end{table}

The low-to-high ratio (50$\times$ on median) is comparable to importance eviction (52.9$\times$), confirming that on long-chain GPQA, the phase gradient itself is largely strategy-independent; what differs is absolute magnitude.

\subsection{Cross-Benchmark Strategy Specificity Summary}

Table~\ref{tab:cross_benchmark} synthesizes the full cross-benchmark picture of strategy specificity.

\begin{table}[h]
\centering
\caption{Cross-benchmark strategy specificity summary (DeepSeek, all at $s=0.30$ except where noted). Strategy specificity (imp/rand ratio) decreases with chain length. On long-chain tasks, phase-dependent severity is near-universal.}
\label{tab:cross_benchmark}
\begin{tabular}{lcccc}
\toprule
Benchmark & Avg.\ Chain & Imp/Rand Ratio ($s$0.20) & Imp/Rand Ratio ($s$0.25) & Imp/Rand Ratio ($s$0.30) \\
\midrule
MATH-500 & $\sim$1200 & $\infty$ (rand NS) & 15.5$\times$ & 16.6$\times$ \\
GSM8K & $\sim$2500 & 2.5$\times$ & 3.8$\times$ & 2.8$\times$ \\
GPQA & $\sim$6000 & 1.52$\times$ & 1.47$\times$ & 1.47$\times$ \\
\bottomrule
\end{tabular}
\end{table}

\section{Discussion}

\paragraph{Universality of phase-dependent severity.}
Our results establish that phase-dependent compound severity is not a property unique to importance-based eviction. Any method that removes tokens from the KV cache produces differential damage across reasoning phases, with late phases (containing synthesis and conclusion steps) consistently more vulnerable. The phenomenon is universal; what varies is magnitude.

\paragraph{Discriminativeness as the scaling variable.}
The importance-to-random ratio provides a natural axis for characterizing eviction methods. On MATH-500, importance scoring achieves high discriminativeness (concentrating eviction on low-utility tokens in early phases while inadvertently removing critical tokens in late phases), yielding a 16.6$\times$ amplification. On GPQA, longer contexts dilute this discriminative advantage because attention distributions become more uniform over many tokens, reducing the ratio to 1.47$\times$. This suggests a scaling law: discriminativeness $\propto$ (context length)$^{-\alpha}$ for some $\alpha > 0$.

\paragraph{Recency as a non-control.}
We initially proposed recency eviction as a scoring-independent negative control. Under closer examination, recency eviction is not independent of reasoning structure: by preserving recently accessed tokens, it preferentially retains the active reasoning frontier. Its null gradient reflects strategy mildness (low effective eviction of critical tokens) rather than absence of phase sensitivity. A more appropriate control would require a strategy that is both non-scoring-based and non-structure-preserving, which random eviction already provides.

\paragraph{Implications for deployment.}
The universality of phase-dependent severity implies that any fixed-rate compound compression configuration will systematically underperform in late reasoning phases. Phase-aware scheduling (reducing compound sparsity during synthesis phases) should benefit all eviction strategies, not only importance-based ones. The magnitude of potential recovery scales with the strategy's discriminativeness on the target workload.

\paragraph{Limitations.}
(1) All models are 7--8B parameters; behavior at larger scales is unknown. (2) The GPQA random $s=0.30$ result is projected from trend extrapolation rather than directly measured. (3) We do not provide a formal discriminativeness metric beyond the empirical imp/rand ratio. (4) No wall-clock throughput measurements are included; this is a measurement study, not a method paper.

\section{Conclusion}

We have demonstrated that phase-dependent compound severity is a universal phenomenon in LLM inference under joint head pruning and KV cache eviction. The magnitude of the phase gradient scales with eviction discriminativeness, decreasing from 16.6$\times$ importance-to-random amplification on medium-chain mathematical reasoning (MATH-500) to 1.47$\times$ on long-chain heterogeneous reasoning (GPQA Diamond). This scaling law, confirmed across three architectures (DeepSeek-R1-Distill-Qwen-7B, Qwen2.5-7B-Instruct, DeepSeek-R1-Distill-Llama-8B), three benchmarks, and four sparsity levels, provides quantitative guidance for compound compression configuration: the longer the reasoning chain, the less discriminative advantage importance scoring provides, and the more uniformly all eviction methods produce phase-dependent damage.

\bibliographystyle{plainnat}
\begin{thebibliography}{10}

\bibitem[Zhang et al.(2023)]{h2o2023}
Zhang, Z., Sheng, Y., Zhou, T., Chen, T., Zheng, L., Cai, R., Song, Z., Tian, Y., Ré, C., Barrett, C., Wang, Z., and Chen, B.
\newblock H2O: Heavy-Hitter Oracle for efficient generative inference of large language models.
\newblock In \emph{NeurIPS}, 2023.

\bibitem[Xiao et al.(2023)]{streamingllm2023}
Xiao, G., Tian, Y., Chen, B., Han, S., and Lewis, M.
\newblock Efficient streaming language models with attention sinks.
\newblock In \emph{ICLR}, 2024.

\bibitem[Xu et al.(2025)]{thinkv2025}
Xu, Y., et al.
\newblock ThinKV: Phase-aware KV cache compression for efficient long-context LLM inference.
\newblock \emph{arXiv:2510.01290}, 2025.

\bibitem[Luo et al.(2025)]{rkv2025}
Luo, Z., et al.
\newblock R-KV: KV cache compression for reasoning-intensive tasks.
\newblock \emph{arXiv:2505.24133}, 2025.

\bibitem[Wang et al.(2024)]{spark2024}
Wang, X., et al.
\newblock SparK: Joint sparsity in attention heads and KV cache for efficient LLM inference.
\newblock In \emph{EMNLP}, 2024.

\bibitem[Chen et al.(2026)]{lethe2026}
Chen, M., et al.
\newblock Lethe: Layer-time adaptive KV cache pruning.
\newblock In \emph{AAAI}, 2026.

\bibitem[Li et al.(2026)]{physicskv2026}
Li, J., et al.
\newblock On the physics of KV cache compression: Phase transitions and representational rigidity.
\newblock \emph{arXiv:2603.01426}, 2026.

\bibitem[Zhao et al.(2025)]{hies2025}
Zhao, R., et al.
\newblock HIES: Entropy-compensated importance scoring for KV cache eviction.
\newblock \emph{arXiv:2510.13832}, 2025.

\end{thebibliography}

\end{document}
